
\Bsec{Kondition und Stabilität} {
		
		\begin{itemize}
			\itemf Fehlerberechnung
			\itemf Kondition
			\itemf Stabilitäts Bestimmung
			\itemf Kombinations Aufgaben und Verständnisfragen
		\end{itemize}
}
%----------------------------------------------------------------------------------
\bsubsec{Fehlerberechnung}{
%		
%		Absoluter Fehler: $$f_{abs} = \abs{x - rd(x)}$$
%		
%		Relativer Fehler: $$f_{rel} = \frac{\abs{x - rd(x)}}{x}$$
%		
%	}
%		
		Rechne den exakten und gerundeten Wert der folgenden Operationen aus:
		
		Zusatz: Bestimme den absoluten und relativen Fehler nach jeder Operation.
		
		Sei 
		\begin{align*}
			a = 2,4142 \quad rd(a) = 2,42 \quad ,
		\end{align*}
		wir haben also 3 signifikante (dezimale) Stellen, runden nach jeder Operation und im Falle dessen immer auf.
		\begin{enumerate}
			\itemf $b = a + 0.2$
			\itemf $c = b + b$
			\itemf $d = c + 0.00123$
			\itemf $e = d - c$
			\itemf Wie nennt man die Probleme bei 3. und 4.?
		\end{enumerate}
%		
	}
	\bsol{
		
		Anfangs Fehler von ca. 0,24\%
		\begin{enumerate}
			\itemf $b = 2,6142; \quad rd(b) = 2,62;$ \quad Fehler von ca. 0,222\%
			\itemf $c = 5,2284; \quad rd(c) = 5,24;$ \quad Fehler von ca. 0,222\%
			\itemf $d = 5,22963; \quad rd(d) = 5,25;$ \quad Fehler von ca. 0,39\%, verursacht durch Absorption
			\itemf $e = 0.00123; \quad rd(e) = 0,01;$ \quad Fehler von ca. 713\% (!), Ursache ist eine Auslöschung
		\end{enumerate}
	}
%----------------------------------------------------------------------------------
\bsubsec{Kondition}{
%		
%		Konditionszahl: $$cond(f, x) = \abs{\frac{x \cdot f'(x)}{f(x)}}$$ \sbreak
%		
%		Eingabefehler: $$\abs{\frac{x - rd(x)}{x}}$$ \sbreak
%			
%		Ausgabefehler: $$\abs{\frac{f(x) - f(rd(x))}{f(x)}}$$ \sbreak
%		
%		Verstärkungsfaktor: $$\frac{rel. Ausgabefehler}{rel. Eingabefehler}$$ \sbreak
%	}
%		
		Rechnen sie für folgenden Aufgaben die Konditionszahl aus und interpretieren sie das Ergebnis:
		\begin{enumerate}
			\itemf $f(x) = 2x^2$
			\itemf $g(x) = 5e^x$
			\itemf $h(x) = \ln(x)$
			\itemf $i(x) = \sqrt{4x} \cdot x$
		\end{enumerate}
		
	}
	\bsol{
		
		\begin{enumerate}
			\itemf $\cond(f, x) = 2$ \\$\Rightarrow$ gut konditioniert
			\itemf $\cond(g, x) = x$ \\$\Rightarrow$ bei $\abs{x} \gg 1$ schlecht konditioniert, ansonsten gut
			\itemf $\cond(h, x) = \frac{1}{\ln(x)}$ \\$\Rightarrow$ schlecht konditioniert, solange $x \le e$ gilt
			\itemf $\cond(i, x) = \frac{x * 3\sqrt{x}}{2 \cdot x^{3/2}} = \nicefrac{3}{2} $ \\$\Rightarrow$ gut konditioniert
		\end{enumerate}
	}

%----------------------------------------------------------------------------------
\bsubsec{Stabilität}{
%		
%		Jede Operation erzeugt einen Fehler:\\
%		$(a$ $op_M$ $b) = (a$ op $b) * (1 + \epsilon)$
%		
%		Endergebnisfehler: $$\abs{\frac{rd(f)(x) - f(x)}{f(x)}}$$
%	}
%		
		Bestimme, ob folgende Funktionen jeweils stabil oder nicht sind: (Verkettungen von Fehlern werden ignoriert)
		\begin{enumerate}
			\itemf $f(x) = 2x^2$
			\itemf $g(x) = 5e^x$ (<- $e^x$ erzeugt \fat{einen} Fehler)
			\itemf $h(x) = \ln(x) \cdot x$
		\end{enumerate}
		
	}
	\bsol{
		
		\begin{enumerate}
			\itemf Endergebnisfehler: $\abs{\epsilon_1 + \epsilon_2} \imp$ Stabil
			\itemf Endergebnisfehler: $\abs{\epsilon_1 + \epsilon_2} \imp$ Stabil
			\itemf Endergebnisfehler: $\abs{\epsilon_1 + \epsilon_2} \imp$ Stabil
		\end{enumerate}
	}
%----------------------------------------------------------------------------------
\bquestion{
%	}
%		
%		Beantworte folgende (Trick-)Fragen möglichst genau:
		\begin{enumerate}
			\itemf Was beschreibt die Konditionszahl?
			\itemf Beschreibe kurz den Unterschied zwischen Kondition und Stabilität.
			\itemf Erläutere, was eine Konditionszahl von $< 1$ bedeuten würde.
			\itemf Aus dem Bauchgefühl: Ist f(x) = x stabil? Warum?
		\end{enumerate}
		
	}
	\bsol{
		
		\begin{enumerate}
			\itemf Die ungünstigste Fehler Verstärkung eines Eingabesfehlers, der bei dem Problem auftreten kann.
			\itemf Kondition: Fehler des Problems, vgl. Eingabe- und Ausgabefehler\\
				Stabilität: Fehler vom Algorithmus, Fehler durch Arithmetik
			\itemf Der Eingabefehler würde sich verkleinern, ist quasi nicht möglich.
			\itemf Keine Operation, also auch keine Fehlererzeugung. Es muss stabil sein!
		\end{enumerate}
	}
	%----------------------------------------------------------------------------------